% 从旧数学附录迁出的筛法与判素说明；不新增算法。
\section{杜教筛：预筛与缓存都不可省略}\label{knowledge-dujiao}\index{Du Jiao sieve}
整除反演建模（第~\ref{knowledge-mobius-model}~节，第~\pageref{knowledge-mobius-model}~页）的两条卷积恒等式给出 \texttt{DuJiao} 的
\texttt{mertens} 与 \texttt{totient\_sum} 递推：商相同的一段合并，系数为区间长度；超过预筛上限的结果必须写回缓存。
最大查询 $N$、预筛上限 $B$ 时单次 $O(B+N/\sqrt B)$，取 $B\approx N^{2/3}$ 得 $O(N^{2/3})$；
多个不同查询会继续增加缓存。范围 $0\le n\le10^{11}$，欧拉函数前缀和用 128 位。
取模时负值先求余，仅余数为负时加模数。

\section{64 位判素与质因数分解}\label{knowledge-miller-bounds}\index{Miller--Rabin}\index{Pollard--Rho}
Miller--Rabin 用七个固定底数，覆盖无符号 64 位整数：
\[
2,\ 325,\ 9375,\ 28178,\ 450775,\ 9780504,\ 1795265022.
\]
$n-1=d2^s$，检查 $a^d$ 及平方链；$n<2$ 非素数；底数模 $n$ 为 0 时跳过。模乘用无符号 128 位。

Pollard--Rho 迭代 $f(x)=x^2+c\pmod n$，取 $\gcd(|x-y|,n)$；只接受 $1<d<n$，否则换参数重试，
无确定最坏上界。\texttt{factor} 返回排序且保留重复的质因子，\texttt{factor(1)} 为空，取首尾前先判空。

\section{单个数的欧拉函数}\label{knowledge-phi-single}\index{Euler totient}
实现见第~\pageref{compact-euler_phi}~页。对 $n\ge1$，
\[
\varphi(n)=n\prod_{p\mid n}\left(1-\frac1p\right),
\]
按 $\mathrm{ans}\gets(\mathrm{ans}/p)(p-1)$ 先除后乘，不会溢出。
试除入口用 $p\le n/p$，$O(\sqrt n)$，不适合 64 位大素数；传入 \texttt{PollardRho} 对象时先分解再去重更新。
$\varphi(1)=1$，0 不在定义域内。

\section{线性筛：三张表的约定}\label{knowledge-linear-sieve}\index{linear sieve}
\texttt{LinearSieve} 构造最小质因子、$\varphi$、$\mu$，每个合数只由 $\operatorname{lp}$ 生成一次，$O(n)$。
若 $p\mid i$，$\varphi(ip)=p\varphi(i)$、$\mu(ip)=0$；否则为 $(p-1)\varphi(i)$ 与 $-\mu(i)$。
判素用 $x\ge2$ 且 $\operatorname{lp}(x)=x$；反复除 $\operatorname{lp}(x)$ 分解，0 不可分解。
要求 $0\le n<\texttt{INT\_MAX}$，乘积前检查 $p\le n/i$。
相关恒等式：$\sum_{d\mid x}\varphi(d)=x$，$\sum_{d\mid x}\mu(d)=[x=1]$，
$\varphi(n)=\sum_{d\mid n}\mu(d)(n/d)$。


\section{区间筛：只保存目标区间}\label{knowledge-segmented-sieve}\index{segmented sieve}\index{区间筛素数}
实现见第~\pageref{compact-segmented_primes}~页。
设 $0\le L\le R\le\texttt{INT\_MAX}$，筛出不超过 $\lfloor\sqrt R\rfloor$ 的素数，对每个 $p$ 从
\[
\max\bigl(p^2,\lceil L/p\rceil p\bigr)
\]
开始标记区间倍数，输出跳过 0 和 1。
令 $W=R-L+1$，空间 $O(\sqrt R+W)$，时间 $O((\sqrt R+W)\log\log(R+3))$；每次调用重建底筛。
$p^2$、取整加法与遍历变量用 \texttt{long long}，否则 $R=\texttt{INT\_MAX}$ 时溢出。
